%!TEX root = ../thesis.tex
% Титульный лист
% \thispagestyle{empty}
\linespread{1.1}

\begin{center}
    {\bfseries
    НАЦІОНАЛЬНИЙ ТЕХНІЧНИЙ УНІВЕРСИТЕТ УКРАЇНИ \\
    <<КИЇВСЬКИЙ ПОЛІТЕХНІЧНИЙ ІНСТИТУТ \\
    імені Ігоря СІКОРСЬКОГО>>\\
    Навчально-науковий фізико-технічний інститут\\
    \kafedra}
\end{center}
\par

\linespread{1.1}
Рівень вищої освіти --- перший (бакалаврський)

Спеціальність --- 113~Прикладна математика,

ОПП <<\op>>

\vspace{10mm}
\begin{tabularx}{\textwidth}{XX}
    & ЗАТВЕРДЖУЮ                               \\[06pt]
    & В.о. завідувача кафедри                 \\[06pt]
    & \rule{2.5cm}{0.25pt} \zavcafPibShort     \\[06pt]
    & <<\rule{0.5cm}{0.25pt}>> \rule{2.5cm}{0.25pt} \YearOfDefence~р. 
\end{tabularx}

\vspace{5mm}
\begin{center}
    {\bfseries ЗАВДАННЯ \\}
    {\bfseries на дипломну роботу \\}
\end{center}

%%%%%====================================
% !!! Не чіпайте наступні три команди!
%%%%%====================================
\frenchspacing
\doublespacing          % інтервал "1,5" між рядками, тепер навічно
\setfontsize{14}

Студент: \underline{\reportAuthor}

\begin{enumerate}[label=\arabic*.]
    \item 
        Тема роботи: <<\emph{\reportTitle}>>, %\\
        науковий керівник: \supervisorRegalia ~\supervisorPib, \par
        затверджені наказом по університету \textnumero \rule{0.5cm}{0.25pt} від <<\rule{0.5cm}{0.25pt}>> \rule{2.5cm}{0.25pt} \YearOfDefence~р.

    \item Термін подання студентом роботи:
        <<\rule{0.5cm}{0.25pt}>> \rule{2.5cm}{0.25pt} \YearOfDefence~р.

    \item Об'єкт дослідження:
        процес інтерактивної багатоваріантної колоризації зображень.

    \item Предмет дослідження:
        моделі глибокого навчання для колоризації зображень,
        метрики для їх порівняльного аналізу,
        а також методи суб'єктивного оцінювання якості інтерактивної колоризації.

    \item Перелік завдань: %\emph{(впишіть теми та задачі, які ви розкриваєте у роботі; можна робити це попунктно)}
        \begin{itemize}[label=$\bullet$]
            \item проаналізувати сучасний стан досліджень у задачі колоризації зображень та виявити обмеження існуючих методів глибокого навчання, зокрема архітектур на базі CNN, GAN, Transformer та Diffusion;
            \item сформувати набори даних для навчання і тестування, а також розробити алгоритм генерації кольорових підказок для імітації взаємодії з користувачем;
            \item спроектувати та реалізувати власну модель колоризації на базі гібридної архітектури U-Net та B-Mamba, а також інтегрувати базові моделі інших архітектур для порівняння;
            \item розробити програмне забезпечення із графічним інтерфейсом для виконання інтерактивної багатоваріантної колоризації;
            \item здійснити навчання розробленої моделі та провести комплексне оцінювання її ефективності шляхом порівняльного тестування за об'єктивними метриками і проведення користувацького дослідження.
        \end{itemize}

    \item 
        Орієнтовний перелік графічного (ілюстративного) матеріалу: Презентація доповіді.

    \item
        Дата видачі завдання: 3 вересня \YearOfBeginning~р.
\end{enumerate}

% Якщо перша частина завдання вилізе за сторінку - приберіть команду \newpage
% Календарний план є продовженням завдання, а не окремою частиною

% \newpage
% \thispagestyle{empty}

\begin{center}
    Календарний план
\end{center}

\renewcommand{\arraystretch}{1.5}
\begin{table}[h!]
    \setfontsize{14pt}
    \centering
    \begin{tabularx}{\textwidth}{|>{\centering\arraybackslash\setlength\hsize{0.25\hsize}}X|>{\centering\setlength\hsize{2\hsize}}X|>{\centering\arraybackslash\setlength\hsize{1\hsize}}X|>{\centering\arraybackslash\setlength\hsize{0.75\hsize}}X|}
        \hline \textnumero \par з/п & Назва етапів виконання дипломної роботи & Термін виконання & Примітка \\
        \hline 
        1 & Узгодження теми роботи із науковим керівником & 01-15 вересня \YearOfBeginning~р. & Виконано \\
        \hline 
        2 & Огляд опублікованих джерел за тематикою дослідження & Вересень-грудень \YearOfBeginning~р. & Виконано \\
        \hline
        3 & Планування кроків проведення дослідження та збирання датасетів & Квітень \YearOfDefence~р. & Виконано \\
        \hline
        4 & Написання програмного забезпечення та проведення досліджень & Квітень-травень \YearOfDefence~р. & Виконано \\
        \hline
        5 & Аналіз отриманих результатів & Травень \YearOfDefence~р. & Виконано \\
        \hline
        6 & Оформлення дипломної роботи & Травень \YearOfDefence~р. & Виконано \\
        \hline
        7 & Попередній зaхист бакалаврської роботи & 2~червня \YearOfDefence~р & Виконано \\
        \hline
        8 & Подача закінченої бакалаврської роботи на кафедру & 4~червня \YearOfDefence~р. & Виконано \\
        \hline
        9 & Захист бакалаврської роботи & 18~червня \YearOfDefence~р. & Виконано \\
        \hline
    \end{tabularx}
\end{table}

\renewcommand{\arraystretch}{1}
\begin{tabularx}{\textwidth}{>{\setlength\hsize{1.2\hsize}}X >{\setlength\hsize{0.5\hsize}}X >{\setlength\hsize{1.3\hsize}}X}
    Студент  & \rule{2.5cm}{0.25pt}  & \reportAuthorShort \\[06pt]
    Керівник & \rule{2.5cm}{0.25pt}  & \supervisorPibShort \\
\end{tabularx}

\newpage
